\input{sol_pre.tex}
\input{q_2026-10-02_2.tex}
\begin{document}
\input{sol_head.tex}

%================= 1
\Savol{1}
\[\overline{abab}=1000a+100b+10a+b=1010a+101b=101(10a+b).\]
Demak son har doim 101 ga bölinadi (11, 111, 100 ga har doim emas: masalan $1212$).
\Javob{B) 101}

%================= 2
\Savol{2}
\[\frac1{\frac12-1}=\frac1{-\frac12}=-2,\qquad \frac1{\frac12+1}=\frac1{\frac32}=\frac23.\]
\[\frac{3-(-2)}{3+\frac23}=\frac5{\frac{11}3}=\frac{15}{11}.\]
\Javob{A) $\frac{15}{11}$}

%================= 3
\Savol{3}
Soat 03:00 da minut mili soat milidan $90^\circ$ orqada. Minut mili minutiga $6^\circ$, soat mili $0{,}5^\circ$ yuradi, ya'ni minut mili har minutda $5{,}5^\circ$ ga yaqinlaşadi:
\[t=\frac{90}{5{,}5}=\frac{180}{11}=16\frac4{11}\ \text{minut}.\]
\Javob{D) $16\frac4{11}$}

%================= 4
\Savol{4}
Abrorning maoşi $x$, Anvarniki $x+500\,000$:
\[0{,}4(x+500\,000)-0{,}3x=350\,000\ \Rightarrow\ 0{,}1x=150\,000\ \Rightarrow\ x=1\,500\,000.\]
Abrorning qöşimça haqi: $0{,}3\cdot1\,500\,000=450\,000$ söm.
\Javob{C) 450\,000 söm}

%================= 5
\Savol{5}
\[a^2b^{-\frac32}\cdot a^{-1}b^{\frac32}:\left(ab^{\frac12}\right)=a\cdot b^0:\left(ab^{\frac12}\right)=b^{-\frac12}=\frac1{\sqrt b}.\]
\Javob{A) $\frac1{\sqrt b}$}

%================= 6
\Savol{6}
$\sqrt{2\pm\sqrt3}=\sqrt{\dfrac{4\pm2\sqrt3}2}=\dfrac{\sqrt3\pm1}{\sqrt2}$. \ Maxrajlar:
\[\sqrt2+\frac{\sqrt3+1}{\sqrt2}=\frac{3+\sqrt3}{\sqrt2},\qquad \sqrt2-\frac{\sqrt3-1}{\sqrt2}=\frac{3-\sqrt3}{\sqrt2}.\]
\[\sqrt2\left(\frac{2+\sqrt3}{3+\sqrt3}+\frac{2-\sqrt3}{3-\sqrt3}\right)=\sqrt2\cdot\frac{(3+\sqrt3)+(3-\sqrt3)}{9-3}=\sqrt2\cdot\frac66=\sqrt2.\]
(Suratlar: $(2+\sqrt3)(3-\sqrt3)=3+\sqrt3$, \ $(2-\sqrt3)(3+\sqrt3)=3-\sqrt3$.)
\Javob{A) $\sqrt2$}

%================= 7
\Savol{7}
\[\sqrt[4]{\sqrt7-\sqrt[4]{24}}\cdot\sqrt[4]{\sqrt7+\sqrt[4]{24}}=\sqrt[4]{7-\sqrt{24}}=\sqrt[4]{\left(\sqrt6-1\right)^2}=\sqrt{\sqrt6-1}.\]
\[\sqrt{\sqrt6-1}\cdot\sqrt{\sqrt6+1}=\sqrt{6-1}=\sqrt5.\]
\Javob{D) $\sqrt5$}

%================= 8
\Savol{8}
Ikkinçi yiğindidan birinçisini ayiramiz: $4d=51-40=11\Rightarrow d=\frac{11}4$.\\
$4a_1+12d=40\Rightarrow 4a_1=40-33=7\Rightarrow a_1=\frac74$. \quad $a_{10}=a_1+9d=\dfrac74+\dfrac{99}4=\dfrac{106}4=26{,}5$.
\Javob{B) 26,5}

%================= 9
\Savol{9}
$3a_2=15\Rightarrow a_2=5$, \ $a_1=5-d$, \ $a_3=5+d$. \ $b_1=5-d$, \ $b_2=6$, \ $b_3=10+d$:
\[36=(5-d)(10+d)\ \Rightarrow\ d^2+5d-14=0\ \Rightarrow\ d=2\ \text{yoki}\ d=-7.\]
$d=2$: $3;\ 6;\ 12$ — ösuvçi. \quad $d=-7$: $12;\ 6;\ 3$ — kamayuvçi. \quad $b_1+b_2+b_3=21$.
\Javob{A) 21}

%================= 10
\Savol{10}
$a^3-b^3-3ab(a-b)=a^3-3a^2b+3ab^2-b^3=(a-b)^3=1{,}1^3=1{,}331$.
\Javob{D) 1,331}

%================= 11
\Savol{11}
Yevklid algoritmi:
\[x^8+x^4+1=x^3\left(x^5+x+1\right)+\left(1-x^3\right),\]
\[x^5+x+1=-x^2\left(1-x^3\right)+\left(x^2+x+1\right),\]
\[1-x^3=(1-x)\left(x^2+x+1\right)+0.\]
Oxirgi noldan farqli qoldiq: $x^2+x+1$.
\Javob{B) $x^2+x+1$}

%================= 12
\Savol{12}
\[\frac{\sin^4\alpha}{\cos^2\alpha}+\frac{\cos^4\alpha}{\sin^2\alpha}+3=\frac{\sin^6\alpha+\cos^6\alpha}{\sin^2\alpha\cos^2\alpha}+3=\frac{1-3\sin^2\alpha\cos^2\alpha}{\sin^2\alpha\cos^2\alpha}+3=\frac1{\sin^2\alpha\cos^2\alpha}=\frac4{\sin^22\alpha}.\]
\Javob{C) $\frac4{\sin^22\alpha}$}

%================= 13
\Savol{13}
$\varphi=\arctg\sqrt7$: \ $\cos\varphi=\dfrac1{\sqrt{1+\tg^2\varphi}}=\dfrac1{\sqrt8}=\dfrac1{2\sqrt2}$.
\[2\sqrt2\sin\left(\frac\pi2+\varphi\right)=2\sqrt2\cos\varphi=2\sqrt2\cdot\frac1{2\sqrt2}=1.\]
\Javob{A) 1}

%================= 14
\Savol{14}
$5\cdot0{,}2^{x-2}=5\cdot5^{2-x}=5^{3-x}=\dfrac{25}{5^{x-1}}$. \ $t=5^{x-1}>0$:
\[t+\frac{25}t=26\ \Rightarrow\ t^2-26t+25=0\ \Rightarrow\ t=1\ \text{yoki}\ t=25\ \Rightarrow\ x=1\ \text{yoki}\ x=3.\]
Yiğindi: $1+3=4$.
\Javob{B) 4}

%================= 15
\Savol{15}
$t=x^{\frac14}\ge0$: \ $t^2-3t+2\le0\Rightarrow (t-1)(t-2)\le0\Rightarrow 1\le t\le2\Rightarrow 1\le x\le16$.\\
Butun yeçimlar: $1,2,\dots,16$ — 16 ta.
\Javob{A) 16}

%================= 16
\Savol{16}
Şart: $x+2\ge0$. \ Kvadratga kötaramiz: $2x+7=x^2+4x+4\Rightarrow x^2+2x-3=0\Rightarrow x=1$ yoki $x=-3$.\\
$x=-3$ da $x+2<0$ — çet ildiz. Yagona ildiz $x=1$.
\Javob{A) 1}

%================= 17
\Savol{17}
$a=\log_23$: \ $\log_224=3+a$, \ $\log_26=1+a$, \ $\log_212=2+a$, \ $\log_23=a$.
\[(3+a)^3-(1+a)^3-(2+a)^3+a^3=(27-1-8)+(27-3-12)a+(9-3-6)a^2=18+12a.\]
$N=\dfrac{18+12a}6=3+2a$, \quad $2^N=2^3\cdot\left(2^{a}\right)^2=8\cdot9=72$.
\Javob{D) 72}

%================= 18
\Savol{18}
$x^2+4x=(x+2)^2-4$, \ $t=|x+2|\ge0$:
\[t^2-4-5t-2=0\ \Rightarrow\ t^2-5t-6=0\ \Rightarrow\ t=6\ \ (t=-1<0).\]
$|x+2|=6\Rightarrow x=4$ yoki $x=-8$. \ Yiğindi: $-4$.
\Javob{A) $-4$}

%================= 19
\Savol{19}
Aniqlaniş sohasi $x>1$. Bu sohada $\log_2(x-1)$ ning işorasi $x-2$ ning işorasi bilan bir xil, şuning uçun
\[(x-3)(x-2)\le0\ \Rightarrow\ 2\le x\le3.\]
Butun yeçimlar $2$ va $3$, yiğindisi $5$.
\Javob{A) 5}

%================= 20
\Savol{20}
$g(x)=\dfrac{f(x)}x=ax^2+bx+c$: \ $g(1)=1$, \ $g(2)=1$, \ $g(3)=\dfrac43$.\\
$g(1)=g(2)$ bölgani uçun $g(x)=1+k(x-1)(x-2)$; \ $g(3)=1+2k=\dfrac43\Rightarrow k=\dfrac16$.
\[g(5)=1+\frac16\cdot4\cdot3=3,\qquad f(5)=5\cdot g(5)=15.\]
\Javob{B) 15}

%================= 21
\Savol{21}
Grafikdan: $x\le a$ qismida $f(0)=1\Rightarrow c=1$.\\
Logarifmik qism $(2;0)$ va $(5;2)$ dan ötadi: $\log_2(2-b)=0\Rightarrow b=1$; tekşiriş: $\log_2(5-1)=2$ — töğri.\\
Vertikal asimptota $x=b=1$ va u $x=a$ ştrix çiziq bilan ustma-ust, ya'ni $a=1$.
\[f(a)=f(1)=|1|+1=2.\]
\Javob{A) 2}

%================= 22
\Savol{22}
$(k^x)'=k^x\ln k$:
\[f'(x)=\frac{8^x\cdot3\ln2+4^x\cdot2\ln2+3\cdot2^x\ln2}{\ln2}=3\cdot8^x+2\cdot4^x+3\cdot2^x.\]
$x_0=\log_23$: $2^{x_0}=3$, \ $f'(x_0)=3\cdot27+2\cdot9+3\cdot3=81+18+9=108$.
\Javob{D) 108}

%================= 23
\Savol{23}
\[\int\frac{2\,dx}{\sqrt{2-x^2}}=2\arcsin\frac x{\sqrt2}+C.\]
$F(-1)=2\arcsin\left(-\dfrac1{\sqrt2}\right)+C=-\dfrac\pi2+C=0\Rightarrow C=\dfrac\pi2$.
\Javob{C) $2\arcsin\frac x{\sqrt2}+\frac\pi2$}

%================= 24
\Savol{24}
\fig{\begin{tikzpicture}[scale=0.42]
\draw[thick] (0,0) rectangle (12,9);
\draw[thick] (3,3) circle (3); \draw[thick] (9,3) circle (3); \draw[thick] (6,7) circle (2);
\fill (3,3) circle (2.5pt) node[below left]{\scriptsize$O_1$}; \fill (6,7) circle (2.5pt) node[above left]{\scriptsize$O$};
\draw[axs] (3,3)--(6,7); \draw[ax] (3,3)--(6,3)--(6,7);
\node[aux,left] at (4.5,5.2){\scriptsize$3+r$}; \node[aux,below] at (4.5,3){\scriptsize$3$}; \node[aux,right] at (6,4.6){\scriptsize$6-r$};
\node[right] at (12,4.5){\small 9}; \node[below] at (6,0){\small 12};
\end{tikzpicture}}
Pastki aylanalar: $4R=12\Rightarrow R=3$, markazi $O_1$ çap tomondan va pastdan 3 masofada.\\
Kiçik aylana radiusi $r$, markazi $O$ örta çiziqda, yuqoridan $r$ masofada: $O$ ning balandligi $9-r$.\\
\textcolor{aux}{Qizil} töğri burçakli uçburçakda: gorizontal katet $6-3=3$, vertikal katet $(9-r)-3=6-r$, gipotenuza $O_1O=3+r$:
\[9+(6-r)^2=(3+r)^2\ \Rightarrow\ 45-12r=9+6r\ \Rightarrow\ r=2.\]
Uzunligi: $2\pi r=4\pi$.
\Javob{B) $4\pi$}

%================= 25
\Savol{25}
$C$ dan içki aylanaga urinmalar uzunligi $CK=CL=p-AB$, bunda $p=\dfrac{8+10+12}2=15$: \ $CK=15-8=7$.\\
$EF$ ham aylanaga urinma, uriniş nuqtasi $M$ bölsin. Bitta nuqtadan ötkazilgan urinmalar teng: $EM=EK$, $FM=FL$. Şuning uçun
\[P_{EFC}=CE+EM+MF+FC=(CE+EK)+(FL+FC)=CK+CL=2\cdot7=14.\]
\Javob{A) 14}

%================= 26
\Savol{26}
Diagonal $\sqrt{8^2+6^2}=10$. Töğri burçakli uçburçakda gipotenuzaning uçidan uriniş nuqtasigaça masofa $p-(\text{qarşi tomon})$, $p=\dfrac{6+8+10}2=12$.\\
Diagonalning pastki çap uçidan: yuqoridagi uçburçakda qarşi tomon $8$ $\Rightarrow$ $12-8=4$; \ pastdagi uçburçakda qarşi tomon $6$ $\Rightarrow$ $12-6=6$.
\[x=6-4=2.\]
\Javob{D) 2}

%================= 27
\Savol{27}
$\triangle AEB$: $EB=\sqrt{13^2-12^2}=5$. \ Törtburçak balandligi — $E$ dan $AB$ gaça masofa: $BC=\dfrac{AE\cdot EB}{AB}=\dfrac{60}{13}$.\\
$\triangle EBC$ ($\angle C=90^\circ$): \ $EC=\sqrt{25-\dfrac{3600}{169}}=\sqrt{\dfrac{625}{169}}=\dfrac{25}{13}$.
\[S=\frac12\cdot BC\cdot EC=\frac12\cdot\frac{60}{13}\cdot\frac{25}{13}=\frac{750}{169}.\]
\Javob{B) $\frac{750}{169}$}

%================= 28
\Savol{28}
Muntazam önburçak taşqi çizilgan aylanani 10 ta teng yoyga böladi: har biri $36^\circ$.\\
$\angle BGC$ — içki çizilgan burçak, $BC$ yoyga tiraladi: $\angle BGC=\dfrac{36^\circ}2=18^\circ$.
\Javob{A) $18^\circ$}

%================= 29
\Savol{29}
$AB$ ning örtasi $M(3;0)$. \ $AB$ ning burçak koeffitsiyenti $\dfrac{-2-2}{5-1}=-1$, perpendikulyarniki $1$.
\[y-0=1\cdot(x-3)\ \Rightarrow\ y=x-3.\]
\Javob{A) $y=x-3$}

%================= 30
\Savol{30}
\fig{\begin{tikzpicture}[scale=0.55]
\draw[thick] (-1,0)--(10,0);
\coordinate (H) at (0,0); \coordinate (A) at (0,5); \coordinate (B) at (5,0); \coordinate (C) at (8.66,0);
\draw[thick] (A)--(B); \draw[thick] (A)--(C); \draw[thick] (A)--(H);
\draw[axs] (B)--(C);
\draw pic[draw,angle radius=5pt]{right angle=A--H--B};
\node[left] at (A){\small$A$}; \node[below] at (H){\small$H$}; \node[below] at (B){\small$B$}; \node[below] at (C){\small$C$};
\node[left] at (0,2.5){\scriptsize 5};
\draw pic[draw,angle radius=14pt,"\tiny$45^\circ$",angle eccentricity=1.9]{angle=A--B--H};
\draw pic[draw,angle radius=20pt,"\tiny$30^\circ$",angle eccentricity=1.6]{angle=A--C--H};
\end{tikzpicture}}
Oğmalarning proyeksiyalari: $HB=AH\ctg45^\circ=5$, \ $HC=AH\ctg30^\circ=5\sqrt3$.\\
$B$ va $C$ tekislikda, $H$ markazli $5$ va $5\sqrt3$ radiusli aylanalarda yotadi. $BC\ge HC-HB$, tenglik $H$, $B$, $C$ bir töğri çiziqda va $B$ oraliqda bölganda:
\[BC_{\min}=5\sqrt3-5=5(\sqrt3-1).\]
\Javob{B) $5(\sqrt3-1)$}

%================= 31
\Savol{31}
Heç qaysi tögarakka bormaydiganlar: $30\cdot0{,}2=6$; \ kamida bittasiga: $24$.\\
Ikkalasiga: $14+18-24=8$. \ Faqat matematikaga: $14-8=6$.
\Javob{C) 6}

%================= 32
\Savol{32}
Bitta kubikda 3 tuşmaslik ehtimoli $\frac56$. Qarama-qarşi hodisa orqali:
\[P=1-\left(\frac56\right)^4=1-\frac{625}{1296}=\frac{671}{1296}.\]
\Javob{D) $\frac{671}{1296}$}

%================= 33-35
\newpage
\Savol{33}
\fig{\begin{tikzpicture}[scale=0.3]
\coordinate (B) at (0,0); \coordinate (C) at (14,0); \coordinate (A) at (5,12); \coordinate (A1) at (5,-12); \coordinate (H) at (5,0);
\draw[thick] (B)--(A)--(C)--(A1)--cycle; \draw[thick,dashed] (B)--(C);
\draw[axs] (A)--(A1);
\coordinate (I) at (6.5,0);
\draw[ax,draw=okgreen] (I) circle (6);
\node[left] at (B){\small$B$}; \node[right] at (C){\small$C$}; \node[above] at (A){\small$A$}; \node[below] at (A1){\small$A'$};
\node[aux] at (4.3,6){\scriptsize 12}; \node[left] at (2.5,6.5){\scriptsize 13}; \node[right] at (9.6,6.5){\scriptsize 15}; \node[below] at (9.5,0){\scriptsize 14};
\end{tikzpicture}}
$AB=13$, $AC=15$, $BC=14$ (aylaniş öqi). Geron: $p=21$, \ $S=\sqrt{21\cdot8\cdot7\cdot6}=84$; \ $BC$ ga balandlik $AH=\dfrac{2\cdot84}{14}=12$.\\
Jism — umumiy asosli ikkita konus: asos radiusi $AH=12$, yasovçilari $13$ va $15$, balandliklari yiğindisi $14$.

\textbf{33.} Töla sirt (ikki konus yon sirtlari): \ $\pi\cdot12\cdot13+\pi\cdot12\cdot15=\pi\cdot12\cdot28=336\pi$.

\textbf{34.} Hajm: \ $\dfrac13\pi\cdot12^2\cdot(BH+HC)=\dfrac13\pi\cdot144\cdot14=672\pi$.

\textbf{35.} Öq kesimi — $ABA'C$ deltoid, tomonlari $13,15,15,13$; $13+15=15+13$, demak unga aylana içki çiziladi va u jismga içki çizilgan şar kesimi böladi (markazi $BC$ öqda).
\[r=\frac{S_{ABA'C}}{p}=\frac{2\cdot84}{28}=6,\qquad V=\frac43\pi\cdot6^3=288\pi.\]
\Javob{33 – B; \ 34 – D; \ 35 – F}

%================= 36
\newpage
\Savol{36}
Simmetriklaş: $t=8-x^2$, unda $9-x^2=t+1$, \ $7-x^2=t-1$:
\[(t+1)^4+(t-1)^4=16\ \Rightarrow\ 2\left(t^4+6t^2+1\right)=16\ \Rightarrow\ t^4+6t^2-7=0.\]
\[\left(t^2+7\right)\left(t^2-1\right)=0\ \Rightarrow\ t^2=1\ \Rightarrow\ t=\pm1.\]
$t=1$: $x^2=7\Rightarrow x=\pm\sqrt7$. \quad $t=-1$: $x^2=9\Rightarrow x=\pm3$.

a) Ildizlar $\pm\sqrt7$, $\pm3$ — 4 ta.

b) Köpaytma: $\sqrt7\cdot(-\sqrt7)\cdot3\cdot(-3)=7\cdot9=63$.
\Javob{a) 4; \quad b) 63}

%================= 37
\newpage
\Savol{37}
Şartlar: $\cos x\ne0$, $\cos3x\ne0$. \ $t=\tg x$, \ $\tg3x=\dfrac{3t-t^3}{1-3t^2}$:
\[37\left(3t-t^3\right)=11t\left(1-3t^2\right)\ \Rightarrow\ t\left(111-37t^2-11+33t^2\right)=0\ \Rightarrow\ t\left(100-4t^2\right)=0.\]
$t=0$ yoki $t^2=25$, ya'ni $\tg x=0$, \ $\tg x=\pm5$. (Bularda $1-3t^2\ne0$, demak $\cos3x\ne0$.)\\
$x=\pi n$, \quad $x=\pm\arctg5+\pi n$, \ $n\in\mathbb Z$.

a) $(-\pi;\pi)$ da: $x=0$; \ $\arctg5$, \ $\arctg5-\pi$; \ $-\arctg5$, \ $\pi-\arctg5$ — jami 5 ta.

b) $[-\pi;\pi]$ da yana $x=\pm\pi$ qöşiladi. Eng kattasi $\pi$, eng kiçigi $-\pi$: \ $\pi+(-\pi)=0$.
\Javob{a) 5; \quad b) 0}

%================= 38
\newpage
\Savol{38}
$x$ ni $\dfrac1x$ ga almaştiramiz:
\[\begin{cases}f(x)+3f\left(\frac1x\right)=4x,\\[2pt] 3f(x)+f\left(\frac1x\right)=\dfrac4x.\end{cases}\]
Ikkinçisini 3 ga köpaytirib, birinçisini ayiramiz: $8f(x)=\dfrac{12}x-4x\Rightarrow f(x)=\dfrac{3-x^2}{2x}$.

a) $f(2)=\dfrac{3-4}{4}=-\dfrac14$.

b) $f(x)=0\Rightarrow 3-x^2=0,\ x\ne0\Rightarrow x=\pm\sqrt3$ — 2 ta ildiz.
\Javob{a) $-\frac14$; \quad b) 2}

%================= 39
\newpage
\Savol{39}
$y=e^{2x-a}\Rightarrow \ln y=2x-a\Rightarrow x=\dfrac{\ln y+a}2$. Demak $f$ va $g$ — özaro teskari funksiyalar, grafiklari $y=x$ töğri çiziqqa nisbatan simmetrik.

a) $a=-3$: $f(x)=e^{2x+3}$. Grafiklar orasidagi eng qisqa masofa $f$ grafigidan $y=x$ gaça eng qisqa masofaning ikki baravari. $f$ grafigidagi nuqtadan $y=x$ gaça masofa
\[d(x)=\frac{e^{2x+3}-x}{\sqrt2}\quad(e^{2x+3}>x).\]
$h(x)=e^{2x+3}-x$, \ $h'(x)=2e^{2x+3}-1=0\Rightarrow e^{2x+3}=\dfrac12\Rightarrow x_0=-\dfrac{3+\ln2}2$ (minimum).
\[h_{\min}=\frac12+\frac{3+\ln2}2=\frac{4+\ln2}2,\qquad \rho=2\cdot\frac{4+\ln2}{2\sqrt2}=\frac{\sqrt2\,(4+\ln2)}2.\]
b) Ösuvçi özaro teskari funksiyalar grafiklari faqat $y=x$ da kesişadi, şuning uçun uriniş nuqtasi $y=x$ da va u yerda urinma $y=x$ ga parallel böladi:
\[e^{2x-a}=x,\qquad 2e^{2x-a}=1\ \Rightarrow\ x=\frac12,\quad e^{1-a}=\frac12\ \Rightarrow\ 1-a=-\ln2\ \Rightarrow\ a=1+\ln2.\]
\Javob{a) $\frac{\sqrt2(4+\ln2)}2$; \quad b) $1+\ln2$}

%================= 40
\newpage
\Savol{40}
Parabola $Ox$ ni $3$ va $-4$ da kesadi: $f(x)=a(x-3)(x+4)$; \ $f(0)=-12a=12\Rightarrow a=-1$:
\[f(x)=-x^2-x+12.\]
$EF$ ning burçak koeffitsiyenti $\dfrac{12-0}{0-(-4)}=3$. Urinmada $f'(x_0)=-2x_0-1=3\Rightarrow x_0=-2$, \ $y_0=f(-2)=10$.

a) $x_0+y_0=-2+10=8$.

b) Urinma: $y=10+3(x+2)=3x+16$, u $Ox$ ni $x=-\dfrac{16}3$ da kesadi. Böyalgan soha urinma, parabola va $Ox$ bilan çegaralangan:
\[S=\underbrace{\frac12\left(-4+\frac{16}3\right)\cdot f_{\text{ur}}(-4)}_{x\in[-\frac{16}3;-4]}+\int_{-4}^{-2}\Bigl((3x+16)-(-x^2-x+12)\Bigr)dx=\frac12\cdot\frac43\cdot4+\int_{-4}^{-2}(x+2)^2dx=\frac83+\frac83=\frac{16}3.\]
\Javob{a) 8; \quad b) $\frac{16}3$}

%================= 41
\newpage
\Savol{41}
$AD=AE+ED=12$. \ $\dfrac{DB}{AB}=\dfrac17\Rightarrow \dfrac{AD}{AB}=\dfrac67\Rightarrow AB=14$, $DB=2$.\\
$\dfrac{AF}{AC}=\dfrac12\Rightarrow AF=FC=3$, $AC=6$.
\[S_{ABC}=\frac12\cdot AB\cdot AC\cdot\sin60^\circ=\frac12\cdot14\cdot6\cdot\frac{\sqrt3}2=21\sqrt3.\]
a) $\triangle BDC$ va $\triangle ABC$ ning $C$ dan tuşirilgan balandligi umumiy:
\[S_{BDC}=\frac{DB}{AB}\cdot S_{ABC}=\frac17\cdot21\sqrt3=3\sqrt3.\]
b) $S_{ADC}=\dfrac67\cdot21\sqrt3=18\sqrt3$, \quad $S_{AEF}=\dfrac12\cdot4\cdot3\cdot\dfrac{\sqrt3}2=3\sqrt3$.
\[S_{CDEF}=S_{ADC}-S_{AEF}=18\sqrt3-3\sqrt3=15\sqrt3.\]
\Javob{a) $3\sqrt3$; \quad b) $15\sqrt3$}

%================= 42
\newpage
\Savol{42}
\fig{\begin{tikzpicture}[scale=0.42]
\draw[thick] (0,-6) rectangle (6,0); \draw[thick] (6,-6) rectangle (12,0); \draw[thick] (3,0) rectangle (9,6);
\coordinate (A) at (0,0); \coordinate (B) at (6,6); \coordinate (C) at (12,-6); \coordinate (P) at (9,0); \coordinate (K) at (6,0);
\draw[thick] (A)--(B)--(C)--cycle;
\draw[axs] (B)--(K); \draw[ax] (A)--(12,0);
\node[left] at (A){\small$A$}; \node[above] at (B){\small$B$}; \node[below right] at (C){\small$C$};
\fill[aux] (P) circle (3pt); \node[aux,above right] at (P){\scriptsize$P$}; \node[aux,below left] at (K){\scriptsize$K$};
\node[aux,right] at (6,3){\scriptsize 6};
\end{tikzpicture}}
$B$ — yuqori kub ustki qirrasining örtasi. Yuqori kub pastkilarning teng yarmida yotgani uçun uning çap qirrasi $A$ dan $3$ cm da.\\
$K$ — $B$ ning $A$ orqali ötuvçi gorizontal çiziqqa proyeksiyasi: $AK=3+3=6$, \ $BK=6$.

a) $\triangle ABK$ — töğri burçakli teng yonli ($AK=BK=6$), şuning uçun $\alpha=45^\circ$.

b) $AK$ çiziq $BC$ ni $P$ da kessin. $B$ bu çiziqdan 6 cm yuqorida, $C$ esa 6 cm pastda, demak $P$ — $BC$ ning örtasi va $B$, $C$ ning proyeksiyalari ($6$ va $12$) örtasida: $AP=9$.\\
$AP$ kesma $\triangle ABC$ ni ikki uçburçakka böladi, ularning $AP$ ga balandliklari $6$ va $6$:
\[S_{ABC}=S_{ABP}+S_{APC}=\frac12\cdot9\cdot6+\frac12\cdot9\cdot6=54.\]
\Javob{a) $45^\circ$; \quad b) 54}

%================= 43
\newpage
\Savol{43}
\fig{\begin{tikzpicture}[scale=0.85]
\coordinate (O) at (0,0);
\coordinate (B) at (90:2); \coordinate (A) at (150:2); \coordinate (E) at (240:2); \coordinate (D) at (330:2); \coordinate (C) at (30:2);
\draw[thick] (O) circle (2);
\draw[thick] (A)--(B)--(C)--(D)--(E)--cycle;
\draw[ax] (B)--(E); \draw[ax] (B)--(D);
\foreach \P in {A,B,C,D,E} \draw[axs,draw=okgreen,thin] (O)--(\P);
\fill (O) circle (1.5pt) node[below right]{\scriptsize$O$};
\node[above] at (B){\small$B$}; \node[left] at (A){\small$A$}; \node[right] at (C){\small$C$}; \node[below] at (E){\small$E$}; \node[right] at (D){\small$D$};
\end{tikzpicture}}
Içki çizilgan burçak özi tiralgan yoyning yarmiga teng:\\
$\angle ABE=45^\circ\Rightarrow \smile AE=90^\circ$; \quad $\angle EBD=45^\circ\Rightarrow \smile ED=90^\circ$.\\
$AB=2=R$ $\Rightarrow$ $\triangle OAB$ teng tomonli, $\smile AB=60^\circ$. \ Qolgan $\smile BCD=360^\circ-90^\circ-90^\circ-60^\circ=120^\circ$, \ $BC=CD\Rightarrow \smile BC=\smile CD=60^\circ$.\\
Vatar: $AE=ED=2R\sin45^\circ=2\sqrt2$, \quad $BC=CD=2R\sin30^\circ=2$.

a) $P=2+2+2+2\sqrt2+2\sqrt2=6+4\sqrt2$.

b) Markaz beşburçak içida (barça yoylar $180^\circ$ dan kiçik). Beşburçakni $O$ markazli 5 ta uçburçakka bölamiz (\textcolor{okgreen}{yaşil}):
\[S=3\cdot\frac12R^2\sin60^\circ+2\cdot\frac12R^2\sin90^\circ=3\cdot2\cdot\frac{\sqrt3}2+2\cdot2=4+3\sqrt3.\]
\Javob{a) $6+4\sqrt2$; \quad b) $4+3\sqrt3$}

%================= 44
\newpage
\Savol{44}
\fig{\begin{tikzpicture}[scale=0.32]
\coordinate (A) at (-4.24,0); \coordinate (B) at (4.24,0); \coordinate (S) at (0,12); \coordinate (O) at (0,0);
\draw[thick] (A)--(S)--(B)--cycle; \draw[ax] (S)--(O);
\draw[thick,draw=okgreen] (0,3) circle (3);
\coordinate (K) at ($(B)!(O)!(S)$);
\draw[axs] (O)--(K);
\draw pic[draw=aux,angle radius=5pt]{right angle=O--K--S};
\node[below] at (O){\small$O$}; \node[left] at (A){\small$A$}; \node[right] at (B){\small$B$}; \node[above] at (S){\small$S$}; \node[right] at (K){\small$K$};
\fill (0,3) circle (2pt) node[left]{\scriptsize$I$};
\end{tikzpicture}}
Asos radiusi $r$, yasovçi $l$: \ $\pi r^2+\pi rl=4\pi r^2\Rightarrow l=3r$, \ $h=\sqrt{9r^2-r^2}=2\sqrt2\,r$.
\[V=\frac13\pi r^2\cdot2\sqrt2\,r=72\pi\ \Rightarrow\ r^3=54\sqrt2=\left(3\sqrt2\right)^3\ \Rightarrow\ r=3\sqrt2,\ \ l=9\sqrt2,\ \ h=12.\]
a) $\triangle SOB$ da $O$ dan gipotenuzaga balandlik $OK$ (\textcolor{aux}{qizil}):
\[OK=\frac{OB\cdot SO}{SB}=\frac{3\sqrt2\cdot12}{9\sqrt2}=4.\]
b) Şar radiusi — öq kesimi $\triangle ASB$ ga içki çizilgan aylana radiusi:
\[\rho=\frac{S_{ASB}}{p}=\frac{3\sqrt2\cdot12}{9\sqrt2+3\sqrt2}=\frac{36\sqrt2}{12\sqrt2}=3,\qquad V_{\text{şar}}=\frac43\pi\cdot27=36\pi.\]
\Javob{a) 4; \quad b) $36\pi$}

%================= 45
\newpage
\Savol{45}
Balandlik $h$, asos radiusi $r$: \ $r^2=4-h^2$, \ $0<h<2$.
\[V(h)=\frac13\pi\left(4-h^2\right)h=\frac\pi3\left(4h-h^3\right),\qquad V'(h)=\frac\pi3\left(4-3h^2\right)=0\ \Rightarrow\ h=\frac2{\sqrt3}.\]
$h<\frac2{\sqrt3}$ da $V'>0$, keyin $V'<0$ — maksimum. \ $r^2=4-\dfrac43=\dfrac83$.

a) $S_{\text{asos}}=\pi r^2=3\cdot\dfrac83=8$.

b) $V_{\max}=\dfrac13\cdot8\cdot\dfrac2{\sqrt3}=\dfrac{16}{3\sqrt3}=\dfrac{16\sqrt3}9$.
\Javob{a) 8; \quad b) $\frac{16\sqrt3}9$}

\ShowJavoblar
\end{document}
